% !TeX root = ../main.tex
% 致谢中主要感谢导师和对论文工作有直接贡献和帮助的人士和单位。
研究生三年时光一晃而过，三年前入学时我激动兴奋的心情至今仍记忆犹新，彼时的我并不清楚研究生生涯究竟会如何度过，此时再来回顾三年的求学生活，生活、科研、求职，千头万绪；迷茫、困顿、憧憬、希望，百感交集；个中滋味，难以尽说，唯有自己清楚。这三年历程，能顺利走到现在，离不开许多人对我的真诚帮助，我要向这些人表示我最诚挚的感谢。

首先要感谢的是我的导师王鸽老师，记得本科阶段有过一个活动，是给三年后的自己写一个时间胶囊，我记得我写过这样一句话：“现在的你如果正在读研究生的话，不知道你遇到了一个怎样的导师？”，现在看来，我很幸运地碰到了一个认真负责，关爱学生的好导师。王老师在科研方面对我的指导极尽认真负责，每周一次的交流讨论风雨无阻，不管是课题大的方向还是实验的细节，王老师都亲自参与讨论与指导。科研之外，王老师在生活中对学生也关爱有加，迷茫时给予指引，焦虑时给予宽慰，消沉时给予鼓励。由衷地感谢王老师在研究生阶段对我提供的所有帮助！

其次要感谢实验室的其他各位老师，丁菡老师和赵衰老师每周的大组会都参与其中，在科研上对我们提出了许多的宝贵意见和指导。惠维老师和赵鲲老师在管理和行政方面做了非常多的工作，确保我们能有一个良好的科研工作环境。同时也要感谢实验室的各位同门和博士师兄，在科研上我们展开了想法的碰撞，在生活中，我们互相帮助，互相扶持，各种有趣的团建活动也点缀了研究生生活。

另外，还要感谢我的父母亲人以及生活中的朋友，他们在背后支撑着我，陪伴着我，让我有更大的勇气和热情来投入到科研工作中来。

最后，感谢我的母校西安交通大学，四年本科时光加三年硕士生涯，这七年里我收获许多，成长许多，保留了许多美好的回忆。同时也感谢各位评审专家、答辩老师和教务工作人员的辛勤付出以及对我提出的宝贵意见。

% 一般致谢的内容有：
% \begin{enumerate}[label=（\chinese*）,itemindent=2em]

%     \item 对指导或协助指导完成论文的导师；
%     \item 对国家科学基金、资助研究工作的奖学金基金、合同单位、资助或支持的企业、组织或个人；
%     \item 对协助完成研究工作和提供便利条件的组织或个人；
%     \item 对在研究工作中提出建议和提供帮助的人；
%     \item 对给予转载和引用权的资料、图片、文献、研究思想和设想的所有者；
%     \item 对其他应感谢的组织和个人。

% \end{enumerate}


\vspace{\baselineskip}

